\chapter{Die offenen Liebesarme}
\quellseite{9}

Der alte Valerius Herberger, der fromme Dichter des Liedes „Valet will ich dir geben“ (No. 426), lag einst auf seinen Knieen am Sterbebett seines sechsjährigen Kindes. Das Ende war nahe, der Vater wußte es, und sein Herz wollte ihm brechen, so lieb hatte er das Kind. Herberger fragte: „Mein Kind, du wirst nun bald von uns gehen und sterben; fürchtest du dich auch vor dem Sterben?“ „Nein, lieber Vater“, antwortete das kranke Kind, „ich fürchte mich nicht zu sterben.“ „Warum ist dir denn nicht bange vor dem Tode?“ fragte der Vater weiter. „Weil ich ja zu Jesu gehe, der mich so lieb hat“, war die treuherzige Antwort. „Aber woher weißt du denn, daß er dich lieb hat?“ fragte der Vater weiter. Und das kranke, todessmatte Kind, das gar nicht verstand, warum der Vater so seltsam fragte, richtete sich mühsam aus seinen Kissen auf, sah den Vater mit großen Augen an, mit so großen, tiefen Augen, wie nur Sterbende sie haben; dann breitete es die Arme aus und sagte in tiefster Bewegung: „So hat er für mich am Kreuze gegangen!“ Und dann sank es erschöpft in die Kissen zurück, und die müden Augen fielen zu, und das müde, kleine Herz stand still. — So hat er auch für mich und dich am Kreuze gehangen. Denkst du daran?
